\clearpage
\renewcommand{\sectioncolor}{bio}
\section{The payoff --- the theorem, assembled \ \ \cell{19}}
\markboth{The assembled theorem}{}

Everything in §5 was preparation for this. The notebook's §9 states and proves the theorem your
Ising notebooks had only checked.

\begin{dobox}[title={Theorem}]
For the 1D Ising transfer matrix $T$ and \textbf{every} $N\ge1$,
\[\operatorname{Tr}(T^N)=\lambda_+^N+\lambda_-^N.\]
\end{dobox}

\begin{center}
\begin{tikzpicture}[font=\footnotesize,
 s/.style={rounded corners=4pt,draw=#1,line width=1.1pt,fill=#1!8,inner sep=8pt,
           text width=7.3cm,align=left,minimum height=1.9cm},
 lic/.style={rounded corners=3pt,fill=#1,text=white,inner sep=5pt,font=\scriptsize\bfseries,
             text width=3.3cm,align=center},
 ar/.style={-{Stealth[length=5pt]},line width=1.1pt,draw=slate!70}]
 \node[s=proof] (s1) at (-3.0,0)     {\textbf{Step 1.} $T$ is real and symmetric.};
 \node[lic=proof] at (5.1,0) {\textbf{direct proof} (§5.1)};
 \node[s=proof] (s2) at (-3.0,-2.35) {\textbf{Step 2.} Hence there is an \emph{orthonormal} basis of
   eigenvectors. Write $P$ for that basis and $D=\operatorname{diag}(\lambda_+,\lambda_-)$, so
   $P^{-1}TP=D$.};
 \node[lic=nano] at (5.1,-2.35) {\textbf{Axler 7.29}\\Real Spectral Theorem};
 \node[s=proof] (s3) at (-3.0,-4.70) {\textbf{Step 3.} $P^{-1}T^NP=(P^{-1}TP)^N=D^N$, and $D^N$ is
   diagonal with entries $\lambda_\pm^N$.};
 \node[lic=sand] at (5.1,-4.70) {\textbf{induction} (§5.6)};
 \node[s=bio] (s4) at (-3.0,-7.05)   {\textbf{Step 4.} The trace is unchanged by a change of basis,
   so $\operatorname{Tr}(T^N)=\operatorname{Tr}(D^N)=\lambda_+^N+\lambda_-^N$. $\blacksquare$};
 \node[lic=bio] at (5.1,-7.05) {\textbf{Axler 10.A}\\trace is basis-free};
 \foreach \a/\b in {s1/s2,s2/s3,s3/s4}{\draw[ar] (\a)--(\b);}
\end{tikzpicture}
\end{center}
\vspace{2pt}
\begin{center}\begin{minipage}{0.93\textwidth}\footnotesize\color{slate}
\textbf{Figure 3.}\; Four steps, each licensed by something named. \textbf{That is what a proof looks
like.} It holds for $N=10^{100}$ with no further work --- which no amount of computation could give
you.
\end{minipage}\end{center}

\begin{warnbox}[title={The gap left deliberately open}]
Step 3 used $(P^{-1}AP)^N=P^{-1}A^NP$ \textbf{without proving it}. That is not an oversight --- it is
notebook \textbf{exercise 6}, and it is your first genuine induction.

Try it before reading on. Base case $N=1$ is immediate; for the step, insert $PP^{-1}=I$ in the
middle and watch it telescope.
\end{warnbox}

\begin{techbox}[title={One theorem, three languages --- \cell{22}, the notebook's §11}]
\begin{center}\small
\begin{tabular}{@{}>{\raggedright\arraybackslash}p{2.4cm}>{\raggedright\arraybackslash}p{8.6cm}>{\raggedright\arraybackslash}p{4.6cm}@{}}
\toprule
\rowcolor{proof!13}
& \textbf{\color{proof}statement} & \textbf{\color{proof}lives in}\\
\midrule
\textbf{Pure} & a self-adjoint operator on a finite-dimensional real inner-product space has an
 orthonormal eigenbasis & Axler 7.29\\
\textbf{Applied} & the free energy is $-\frac1\beta\ln\lambda_+$ and $\xi=1/\ln(\lambda_+/\lambda_-)$ &
 Feynman ch.~5; PBoC2 ch.~6\\
\textbf{Computational} & \bk{T0.eigenmatrix\_right()} returns $(D,P)$ with $P$ orthogonal & Sage\\
\bottomrule
\end{tabular}
\end{center}
\vspace{3pt}
The \textbf{pure} statement licenses the applied one. The \textbf{computational} one only checks
instances. Losing track of which you are doing is the commonest error in computational science.
\end{techbox}
